\documentclass{article}
\usepackage[T1]{fontenc}
\usepackage{lmodern}
\usepackage{graphicx}
\usepackage{amsmath,amssymb}
\usepackage[a4paper,margin=2cm]{geometry}
\usepackage[absolute,overlay]{textpos}
\setlength{\TPHorizModule}{1mm}
\setlength{\TPVertModule}{1mm}
\usepackage[indonesian]{babel}
\usepackage{hyperref}
\setlength{\emergencystretch}{2em}
% unit: Halaman 44

\begin{document}

% === Halaman 44 (python/native) ===

• Di dalam setiap byte bit-bitnya tersusun dari kiri ke kanan dalam urutan yang 

kurang berarti (least significant bits atau LSB) hingga bit-bit yang berarti (most 

significant bits atau MSB). 

• Susunan bit pada setiap byte adalah b7b6b5b4b3b2b1b0. 

 Contoh: 



11010001




  MSB  

       LSB

44

LSB = Least Significant Bit

MSB = Most Significant Bit

Nilai dalam desimal = 1 x 27 + 1 x 26 + 0 x 25 + 1 x 24 + 0 x 23 + 0 x 22 + 0 x 21 + 1 x 20  = 209 

Metode LSB

\end{document}
